%% IIT Mandi lab report template. Fill in the experiment and student details,
%% then write each section. Delete the sections your course does not need.
\documentclass[11pt,a4paper]{article}

\usepackage[a4paper,margin=25mm]{geometry}
\usepackage[T1]{fontenc}
\usepackage{lmodern}
\usepackage{amsmath,amssymb}
\usepackage{siunitx}
\usepackage{graphicx}
\usepackage{xcolor}
\usepackage{booktabs}
\usepackage{caption}
\usepackage[hidelinks]{hyperref}

%% ---------------------------------------------------------------- details
\newcommand{\coursecode}{IC 161P}
\newcommand{\coursename}{Applied Electronics Lab}
\newcommand{\experimentnumber}{3}
\newcommand{\experimenttitle}{Frequency Response of an RC Low-Pass Filter}
\newcommand{\studentname}{Your Name}
\newcommand{\rollnumber}{B25001}
\newcommand{\labgroup}{Group 4}
\newcommand{\instructor}{Dr.\ Instructor Name}
\newcommand{\experimentdate}{12 August 2026}

\definecolor{iitmnavy}{HTML}{00237D}
\setlength{\parindent}{0pt}
\setlength{\parskip}{5pt}

\begin{document}

%% ---------------------------------------------------------------- heading
\begin{center}
  \includegraphics[width=30mm]{iitmandi-logo}\par
  \medskip
  {\small\textsc{Indian Institute of Technology Mandi}\par}
  {\small\coursecode: \coursename\par}
  \bigskip
  {\Large\bfseries\color{iitmnavy} Experiment \experimentnumber\par}
  \smallskip
  {\Large\bfseries\experimenttitle\par}
\end{center}

\medskip
\begin{tabular}{@{}ll@{\hspace{2cm}}ll@{}}
  \textbf{Name:} & \studentname & \textbf{Roll number:} & \rollnumber \\
  \textbf{Group:} & \labgroup & \textbf{Date:} & \experimentdate \\
  \textbf{Instructor:} & \instructor & & \\
\end{tabular}

\noindent\rule{\textwidth}{0.6pt}

%% ---------------------------------------------------------------- report
\section{Aim}
To measure the frequency response of an RC low-pass filter and find its
cut-off frequency.

\section{Apparatus}
Function generator, oscilloscope, resistor (\SI{1}{\kilo\ohm}), capacitor
(\SI{100}{\nano\farad}), breadboard and connecting wires.

\section{Theory}
The cut-off frequency of an RC low-pass filter is
\begin{equation}
  f_c = \frac{1}{2\pi R C} .
  \label{eq:cutoff}
\end{equation}

\section{Procedure}
\begin{enumerate}
  \item Build the circuit and connect the input and output to the oscilloscope.
  \item Vary the input frequency and record the output amplitude.
\end{enumerate}

\section{Observations}
\begin{table}[htbp]
  \centering
  \caption{Output amplitude for an input of \SI{2}{\volt} peak-to-peak.}
  \label{tab:observations}
  \begin{tabular}{@{}S[table-format=5.0]S[table-format=1.2]S[table-format=1.3]@{}}
    \toprule
    {Frequency (\si{\hertz})} & {Output (\si{\volt})} & {Gain} \\
    \midrule
    100   & 1.99 & 0.995 \\
    1000  & 1.70 & 0.850 \\
    10000 & 0.32 & 0.160 \\
    \bottomrule
  \end{tabular}
\end{table}

\section{Calculations}
From Equation~\eqref{eq:cutoff}, $f_c = 1/(2\pi \times \SI{1}{\kilo\ohm}
\times \SI{100}{\nano\farad}) \approx \SI{1.59}{\kilo\hertz}$.

\section{Results}
The measured cut-off frequency was \SI{1.5}{\kilo\hertz}, within 6\% of the
calculated value.

\section{Discussion}
Explain the sources of error, the precautions taken and what you learned.

\end{document}
